\subsection{El punto estacionario del hijo}\label{sec:B4T02}
Conservamos la notación de la Sección~\ref{sec:B4T01}. Ahora $B$ es la matriz RC línea--casilla del
hijo y $D$ es su incidencia casilla--diagonal. La Sección~\ref{sec:B3T07} construye
\[
 \bar p=w+z,\quad B\bar p=\mathbf1,\quad
 a=z/w\in\operatorname{im}B^T,\quad \|a\|_\infty=O(k^{-1}).
\]
Sean $d$ el vector de pérdidas en las diagonales retenidas y $u=d-D^Tz$.
Las vacancias de $\bar p$ son $h+u\ge\eta/2$.

\paragraph{Residuo de estacionariedad tras el reequilibrado.}
Por la estacionariedad del padre y las expresiones exactas de los logaritmos,
\[
 \nabla H_{\rm hijo}(\bar p)\equiv
 D\log(1+u/h)+a-\log(1+a)-3\log(1-z/(1-w))
 \pmod{\operatorname{im}B^T}.
\]
En particular existe un representante $e$ del gradiente tangente tal que
\[
 e=D(u/h)+O_\infty(k^{-2}),\qquad \|e\|_\infty=O(k^{-1}).
\]
Corregir las masas RC no anula en general el primer término diagonal.
No se afirma que $p^*-\bar p=O_{\rm rel}(k^{-2})$.

\paragraph{Paso de Newton y control uniforme.}
Sean $Z$ una base ortonormal de $\ker B$ y
\[
 K=\operatorname{diag}(1/\bar p-3/(1-\bar p))
       +D\operatorname{diag}((h+u)^{-1})D^T,\quad
 R=Z(Z^TKZ)^{-1}Z^T,\quad s=Re.
\]
Aplicamos la estimación matricial de la Sección~\ref{sec:B3T05} al punto factible $\bar p$ con la
brecha heredada. El hijo tiene $k-1\ge\delta_0 n/2$ filas para $k\ge2$, por lo que podemos usar
$\delta_0/2$ incluso en el extremo $k=\delta_0n$. Así obtenemos
$\|\operatorname{diag}(\bar p)^{-1}R\|_{\infty\to\infty}\le M$.
Denotamos $\|v\|_{\bar p}=\|v/\bar p\|_\infty$ y
$\sigma=\|s\|_{\bar p}=O(k^{-1})$.

Para $v\in\ker B$ y $\|v\|_{\bar p}\le r$ suficientemente pequeño,
\[
 \nabla H(\bar p+v)=\nabla H(\bar p)-Kv+\mathcal N(v),
\]
\[
 \|\mathcal N(v)\|_\infty\le Cr^2,\quad
 \|\mathcal N(v)-\mathcal N(w)\|_\infty\le Cr\|v-w\|_{\bar p}.
\]
Las constantes no dependen de una cota inferior de $\bar p$:
$|(D^Tv)_d|\le r\sum_{e\in d}\bar p_e\le r$, y el logaritmo de cada
casilla se desarrolla en $v_e/\bar p_e$, no en $v_e$ solo.
La estacionariedad es equivalente a $v=s+R\mathcal N(v)$.
Para $k$ grande, el segundo miembro define una aplicación que preserva la bola de radio $2\sigma$
y tiene constante de Lipschitz menor que uno; si $\sigma=0$, $\bar p$ ya es estacionario.
Existe por ello un punto $p^*$ con
\[
 \boxed{\|(p^*-\bar p)/\bar p\|_\infty=O(k^{-1}),\qquad
 \|(p^*-\bar p-s)/\bar p\|_\infty=O(k^{-2}).}
\]
Este punto permanece factible y positivo, con vacancias uniformes y $p^*<1/4$.
La hessiana cambiada de signo es definida positiva en la región convexa $p<1/4$,
así que el punto estacionario es único en esa rama. No se afirma unicidad fuera de ella.
El punto $\widehat p=\bar p+s$ ya aproxima al hijo con error relativo $O(k^{-2})$.
Todos los controles son uniformes en los hijos de un padre de la clase.

\paragraph{Incremento de la acción desde $\bar p$.}
En $\bar p$, el término lineal $e^Ts=s^TKs$ no se anula en general. Se combina con
el término cuadrático negativo de la entropía:
\[
 \boxed{H(p^*)-H(\bar p)=\tfrac12 e^TRe+O(k^{-2}).}
\]
El término principal es $O(k^{-1})$. El resto cúbico es
$O(k\|v\|_{\bar p}^3)=O(k^{-2})$, pues $\sum\bar p=k-1$
y hay $O(k)$ diagonales. Los errores al reemplazar $v$ por $s$ también son
$O(k^{-2})$; las partes normales RC se anulan contra todos los incrementos.
